% !TeX root = ../../Book1.tex
\chapter{拓扑与度量空间基础}
\label{b1:app:A}

\begin{chapterabstract}
本附录收集全卷使用的拓扑基础，重点不在列举公理，而在解释几种常用构造：用 Cauchy 列补齐空间，用有限球网辨认紧性，用可数基控制开集，用闭集分离构造连续函数。最后把这些工具结合成波兰空间的稳定性，说明它为什么适合测度收敛问题。
\end{chapterabstract}

\begin{learninggoals}
\begin{itemize}
\item\label{b1:item:A1-goal-metric}区分拓扑性质与特定度量的完备性，掌握完备化及其延拓性质。
\item\label{b1:item:A1-goal-compact}能以全有界性和完备性证明紧性，知道局部紧、相对紧与有界的区别。
\item\label{b1:item:A1-goal-countable}理解可分性、第二可数性和乘积拓扑中各自的可数化步骤。
\item\label{b1:item:A1-goal-extension}能从 Urysohn 引理推出 Tietze 延拓，并为波兰空间的开子集构造相容完备度量。
\end{itemize}
\end{learninggoals}

本附录可在第\ref{b1:ch:01}章之前作为基础补读，也可按正文需要查阅。度量与紧性部分只需实数完备性和集合运算；乘积部分回顾第\ref{b1:sec:0903}节已经证明的 Tychonoff 定理，不把其证明转嫁到本附录。最后一节用到的 Baire 纲定理可在第\ref{b1:sec:0803}节找到。一般点集拓扑可继续读熊金城《点集拓扑讲义》第五版\cite[第3、5--7章]{Xiong2020Topology}；与分析的联系可对照 Hunter、Nachtergaele《Applied Analysis》\cite[第1、2、4章]{Hunter2001Applied}，测度论常用的拓扑事实可查 Fremlin《Measure Theory》\cite[卷四，4A2]{Fremlin2016Measure}。

% !TeX root = ../../Book1.tex
\section{距离、拓扑与完备化}
\label{b1:sec:A101}

分析中经常先给出一种距离，再谈极限、连续与闭集。但这些概念对距离的依赖并不相同：连续性只看开集，Cauchy 条件却比较尚未知道极限的两个点。本节把这一区别落实到构造中。以下用 \(d\) 表示距离；正文中同一个字母常表示欧氏空间的维数，两种用法由所讨论的空间区分。

\subsection{开集与连续性}
\label{b1:subsec:A10101}

\begin{definition}\label{b1:def:app-metric}
设 \(X\) 为集合。若函数 \(d:X\times X\rightarrow[0,\infty)\) 满足
\[
 d(x,y)=0\ \Longleftrightarrow\ x=y,\quad
 d(x,y)=d(y,x),\quad d(x,z)\le d(x,y)+d(y,z),
\]
则称 \(d\) 为 \(X\) 上的\term[duliang]{度量}[metric]，称 \((X,d)\) 为度量空间。开球记作 \(B_d(x,r)=\{\,y:d(x,y)<r\,\}\)。若 \(U\subset X\) 的每个点都包含在一个仍位于 \(U\) 内的开球中，则称 \(U\) 为开集。
\end{definition}

赋范空间中的 \(d(x,y)=\|x-y\|\) 是例子，任意集合上的离散距离也是例子：不同点的距离规定为 \(1\)。反过来，度量空间不必允许向量相加。讨论曲线、测度的状态空间时，这一点尤其有用。

开集对任意并、有限交封闭，并包含 \(\varnothing\) 与 \(X\)。有限交的断言只需在交集中的一点取有限个半径的最小值；对无穷交，这个最小值可能变成零。例如 \(\bigcap_n(-1/n,1/n)=\{0\}\) 不是实数中的开集。

\begin{definition}\label{b1:def:app-topology}
集合 \(X\) 的一个子集族 \(\mathcal T\)，若包含 \(\varnothing,X\)，并对任意并与有限交封闭，则称为 \(X\) 上的\term[tuopu]{拓扑}[topology]。其成员称为开集，开集的补集称为闭集。若集合 \(V\) 包含点 \(x\) 的一个开集邻域，则称 \(V\) 为 \(x\) 的邻域。包含 \(E\) 的全部闭集之交是其闭包 \(\overline E\)；包含于 \(E\) 的全部开集之并是其内部 \(E^\circ\)。
\end{definition}

子集 \(A\subset X\) 自带的子空间拓扑，由所有 \(A\cap U\)、\(U\in\mathcal T\) 组成。因此“闭”必须说明相对于谁。例如 \([1/2,1)\) 在 \((0,1)\) 中闭，却不在 \(\mathbb R\) 中闭。若 \(X\) 来自度量，子空间拓扑正是限制距离所得的拓扑。

\begin{proposition}\label{b1:prop:app-continuity}
映射 \(f:X\rightarrow Y\) 连续，当且仅当每个开集 \(V\subset Y\) 的反像 \(f^{-1}(V)\) 开。这里连续指：对每个 \(x\) 及 \(f(x)\) 的邻域 \(V\)，存在 \(x\) 的邻域 \(U\)，使 \(f(U)\subset V\)。当源与目标均为度量空间时，这又等价于保持所有收敛序列的极限。
\end{proposition}
\begin{proof}
若反像开，则 \(f^{-1}(V^\circ)\) 给出所需邻域。反之，若 \(V\) 开，对反像中的每个点可取映入 \(V\) 的开邻域，反像是这些邻域的并。

连续映射保持序列极限，直接由邻域定义得到。若度量空间间的映射在 \(x\) 不连续，则存在 \(\varepsilon>0\)，对每个 \(n\) 都能取 \(x_n\in B_d(x,1/n)\)，而目标距离满足 \(d_Y\bigl(f(x_n),f(x)\bigr)\ge\varepsilon\)。这个序列反驳保持极限的性质。
\end{proof}

最后一步依靠可数个球控制一点附近的全部邻域。一般拓扑空间不能不加条件地以序列代替连续性的定义；第\ref{b1:sec:0901}节引入网，正是为了处理这种局限。

\subsection{相同拓扑与不同的 Cauchy 列}
\label{b1:subsec:A10102}

\begin{definition}\label{b1:def:app-complete}
若对每个 \(\varepsilon>0\)，存在 \(N\)，使 \(m,n\ge N\) 时 \(d(x_m,x_n)<\varepsilon\)，则称 \((x_n)\) 为 Cauchy 列。若空间中每个 Cauchy 列都收敛到本空间的点，则称它\term[wanbei]{完备}[complete]。
若对每个 \(\varepsilon>0\) 可取不依赖于点的 \(\delta>0\)，使 \(d_X(x,y)<\delta\) 蕴含 \(d_Y\bigl(f(x),f(y)\bigr)<\varepsilon\)，则称 \(f\) 一致连续。
\end{definition}

度量 \(d\) 与 \(\min(1,d)\) 给出相同的开集，也有相同的 Cauchy 列，因为小于 \(1\) 的距离阈值没有改变。相同拓扑本身却不足以推出后一结论。在 \((0,1)\) 上，通常距离与
\[
 \rho(x,y)=\left|\log\frac{x}{1-x}-\log\frac{y}{1-y}\right|
\]
给出同一拓扑；映射 \(x\mapsto\log\frac{x}{1-x}\) 把后一空间等距送到 \(\mathbb R\)，所以 \(\rho\) 完备。通常距离下的 \(1/(n+1)\) 是 Cauchy 列，在 \(\rho\) 下不是。故完备性首先是选定距离的性质。

\begin{proposition}\label{b1:prop:app-complete-closed}
度量空间的完备子空间必为闭集；完备度量空间的闭子空间仍完备。一个 Cauchy 列只要有收敛子列，整个序列便收敛于同一点。
\end{proposition}
\begin{proof}
对最后一句，取子列中的一项同时足够接近极限，又足够深入 Cauchy 尾部，以三角不等式控制全部尾项。
若 \(A\) 完备而 \(x\in\overline A\)，从 \(A\cap B_d(x,1/n)\) 取 \(a_n\)。它是 Cauchy 列，在 \(A\) 中有极限；极限唯一，故 \(x\in A\)。反之，闭子集中的 Cauchy 列先在全空间收敛，闭性再把极限留在子集中。
\end{proof}

“闭子空间完备”不能省掉全空间完备。例如 \(\mathbb Q\) 在自身中既开又闭，却有趋向 \(\sqrt2\) 的 Cauchy 列而没有相应极限。完备化就是把这种缺失的极限系统地补入。

\subsection{由 Cauchy 列补入极限}
\label{b1:subsec:A10103}

\begin{theorem}[度量空间的完备化]\label{b1:thm:app-completion}
每个度量空间 \(X\) 都可等距嵌入一个完备度量空间 \(\widehat X\)，并在其中稠密。这样的空间在保持 \(X\) 中各点的等距同构意义下唯一。
\end{theorem}
\begin{proof}
空空间无需补点。以下设 \(X\ne\varnothing\)，令 \(\mathcal C\) 为 \(X\) 中所有 Cauchy 列组成的集合。在 \(\mathcal C\) 上规定
\[
 (x_n)\sim(y_n)\quad\Longleftrightarrow\quad d(x_n,y_n)\longrightarrow0.
\]
三角不等式给出传递性，因此这确为等价关系。令 \(\widehat X=\mathcal C/{\sim}\)，定义
\[
 \widehat d\bigl([(x_n)],[(y_n)]\bigr)=\lim_{n\to\infty}d(x_n,y_n).
\]
由
\[
 |d(x_n,y_n)-d(x_m,y_m)|
 \le d(x_n,x_m)+d(y_n,y_m)
\]
可知右侧极限存在。再用同一估计比较两组代表，极限不随代表改变。对称性与三角不等式从逐项距离传到极限，零距离的情形恰好是同一个等价类。于是 \(\widehat d\) 是度量。

把 \(x\) 送到常值列的等价类，得到等距嵌入 \(i:X\rightarrow\widehat X\)。若 \(\xi=[(x_n)]\)，则
\[
 \widehat d\bigl(i(x_m),\xi\bigr)=\lim_{n\to\infty}d(x_m,x_n)
 \longrightarrow0\quad(m\to\infty).
\]
所以 \(i(X)\) 稠密。

为证完备，设 \((\xi_m)\) 是 \(\widehat X\) 中的 Cauchy 列。由稠密性取 \(a_m\in X\)，使 \(\widehat d\bigl(i(a_m),\xi_m\bigr)<1/m\)。三角不等式说明 \((a_m)\) 是 \(X\) 中的 Cauchy 列。令 \(\xi=[(a_m)]\)，上段已证 \(i(a_m)\rightarrow\xi\)，故 \(\xi_m\rightarrow\xi\)。

最后设 \(j:X\rightarrow Y\) 是另一稠密等距嵌入，且 \(Y\) 完备。规定
\[
 U([(x_n)])=\lim_{n\to\infty}j(x_n).
\]
完备性保证极限存在，等距性保证良定义并保持距离。对任意 \(y\in Y\)，用 \(j(X)\) 中的序列逼近 \(y\)，可知 \(U\) 满射。它在稠密子集上满足 \(U\circ i=j\)，故也被这一条件唯一确定。
\end{proof}

这称为\term[wanbeihua]{完备化}[completion]。构造并没有给每个 Cauchy 列增加一个不同的新点：彼此距离趋零的列必须识别，否则同一个极限会被重复加入。实数可由有理数作这样的构造；函数空间中用范数完成时，背后的距离步骤也相同，额外的线性运算还需验证能传到等价类。

\begin{proposition}[一致连续延拓]\label{b1:prop:app-uniform-extension}
设 \(A\) 在度量空间 \(X\) 中稠密，\(Y\) 完备。每个一致连续映射 \(f:A\rightarrow Y\) 都有唯一的一致连续延拓 \(F:X\rightarrow Y\)。若 \(f\) 是 \(L\)-Lipschitz 的，则 \(F\) 保持相同常数。
\end{proposition}
\begin{proof}
对 \(x\in X\)，取 \(a_n\in A\) 趋向 \(x\)。一致连续性将 Cauchy 列送到 Cauchy 列，故可令 \(F(x)=\lim_n f(a_n)\)。两列交错后仍是 Cauchy 列，说明定义与逼近列无关。它显然延拓 \(f\)，连续延拓的唯一性也由稠密性得到。

给定 \(\varepsilon>0\)，以 \(\varepsilon/2\) 在 \(A\) 上的一致连续性选取 \(\delta>0\)。若 \(d_X(x,y)<\delta/3\)，可分别用 \(a_n,b_n\) 逼近它们，尾部满足 \(d_X(a_n,b_n)<\delta\)。传到极限得 \(d_Y(F(x),F(y))\le\varepsilon/2<\varepsilon\)。Lipschitz 断言则把原不等式直接传到极限。
\end{proof}

单纯连续不够：\(x\mapsto1/x\) 在稠密子集 \((0,1]\subset[0,1]\) 上连续，却不能延拓到端点。这里的障碍是逼近端点时没有有限极限，与稍后从闭集延拓的 Tietze 定理是不同问题。

% !TeX root = ../../Book1.tex
\section{紧性、有限逼近与局部紧性}
\label{b1:sec:A102}

紧性把无穷的信息压缩成有限的信息。开覆盖的定义最适合一般拓扑，度量空间中的有限球网则把它变成一种可以构造的逼近。本节说明这两种表述怎样联系，也说明只控制集合的直径为什么远远不够。

\subsection{开覆盖与有限球网}
\label{b1:subsec:A10201}

\begin{definition}\label{b1:def:app-compact}
若拓扑空间 \(K\) 的每个开覆盖都有有限子覆盖，则称 \(K\) 紧。对 \(A\subset X\)，紧性按子空间拓扑理解。若度量空间 \(A\) 对每个 \(\varepsilon>0\) 都能被有限个以 \(A\) 中的点为中心、半径为 \(\varepsilon\) 的球覆盖，则称 \(A\) \term[quanyoujie]{全有界}[totally bounded]，这些中心组成一个有限 \(\varepsilon\)-网。
\end{definition}

一个有限 \(\varepsilon\)-网不必逼近所有尺度；要求是每个正尺度都有一个有限网，其点数可以随尺度缩小而增加。例如区间中的等分点可以给出有限网，而无穷离散空间在半径小于 \(1\) 时必须逐点覆盖。后者有界，却不全有界。

\begin{lemma}\label{b1:lem:app-compact-properties}
紧空间的闭子集紧，紧空间在连续映射下的像紧；Hausdorff 空间的紧子集闭。连续双射从紧空间到 Hausdorff 空间时，逆映射也连续。
\end{lemma}
\begin{proof}
对闭子集的开覆盖，加入其补集便覆盖全空间；取有限子覆盖后删去补集即可。对连续像的开覆盖取反像，便归结为定义。

设 \(K\subset X\) 紧而 \(x\notin K\)。对每个 \(y\in K\)，用 Hausdorff 性取互不相交的开邻域 \(U_y\ni y\)、\(V_y\ni x\)。有限个 \(U_{y_j}\) 覆盖 \(K\)，而 \(\bigcap_jV_{y_j}\) 是 \(x\) 的开邻域且不交 \(K\)。所以 \(K\) 闭。最后，连续双射把源空间的闭集送到目标中的紧集，因而送到闭集，恰好说明逆映射连续。
\end{proof}

\begin{theorem}[度量空间的紧性判据]\label{b1:thm:app-metric-compact}
对度量空间 \(X\)，下列条件等价：
\begin{equivalences}
\item\label{b1:item:app-compact-cover}\(X\) 紧；
\item\label{b1:item:app-compact-sequence}每个序列都有在 \(X\) 中收敛的子列；
\item\label{b1:item:app-compact-complete} \(X\) 完备且全有界。
\end{equivalences}
\end{theorem}
\begin{proof}
先由\itemref{b1:item:app-compact-cover}推出\itemref{b1:item:app-compact-sequence}。给定序列，令 \(F_n=\overline{\{x_j:j\ge n\}}\)。这些非空闭集递减；若交为空，补集构成没有有限子覆盖的开覆盖。因此有 \(x\in\bigcap_nF_n\)。递归选取严格递增的 \(j_k\)，使 \(d(x_{j_k},x)<1/k\)。

由\itemref{b1:item:app-compact-sequence}，Cauchy 列有收敛子列，于是整列收敛。若不全有界，则某个 \(\varepsilon>0\) 下不存在有限网。依次在前面各点的 \(\varepsilon\)-球之外选新点，得到两两距离至少为 \(\varepsilon\) 的序列，与存在收敛子列矛盾。这得到\itemref{b1:item:app-compact-complete}。

现在假设\itemref{b1:item:app-compact-complete}。每个序列都有 Cauchy 子列：先取覆盖 \(X\) 的有限个半径 \(1/2\) 球，选一个含无穷多项的球；在这些项中再用有限个半径 \(1/4\) 球选择，依次进行。每次保留的指标集包含于上一次，取对角子列，其足够后的两项同处一个半径 \(2^{-k}\) 的球。完备性给出收敛子列。

还需从子列性质回到开覆盖。给定开覆盖 \(\mathcal U\)，若对每个 \(n\) 都存在 \(x_n\)，使球 \(B_d(x_n,1/n)\) 不包含在任何一个成员中，取子列 \(x_{n_j}\rightarrow x\) 并取 \(U\in\mathcal U\)、\(B_d(x,r)\subset U\)。当 \(j\) 足够大，\(d(x_{n_j},x)<r/2\) 且 \(1/n_j<r/2\)，对应小球便包含于 \(U\)，矛盾。因此存在 \(\delta>0\)，使每个半径为 \(\delta\) 的球均位于某个覆盖成员内。再由全有界性，用有限个半径 \(\delta/2\) 的球覆盖 \(X\)，为各中心选择包含其 \(\delta\)-球的成员，就得到有限子覆盖。
\end{proof}

证明的最后一步常称为 Lebesgue 数论证。它不是说开覆盖中某个成员含有整个空间，而是说每个足够小的球都能被某个成员容纳，允许成员依球心变化。

\begin{corollary}\label{b1:cor:app-relative-compact}
完备度量空间中的子集 \(A\) 相对紧，当且仅当它全有界。
\end{corollary}
\begin{proof}
若 \(\overline A\) 紧，则有限小球覆盖限制到 \(A\) 后，把与 \(A\) 相交的球心移到交点，半径至多扩大一倍，便得全有界性。反之，\(A\) 的有限 \(\varepsilon/2\)-网给出 \(\overline A\) 的有限 \(\varepsilon\)-网；闭包又完备，故紧。
\end{proof}

例如 \((0,1)\subset\mathbb R\) 相对紧但不完备，所缺极限被其闭包补上。在非完备的环境 \(\mathbb Q\) 中，\(\mathbb Q\cap[0,1]\) 全有界，其在 \(\mathbb Q\) 内的闭包却不紧。因此相对紧性仍与所在环境有关。

\subsection{紧性提供的一致控制}
\label{b1:subsec:A10202}

\begin{proposition}[Heine--Cantor]\label{b1:prop:app-heine-cantor}
从紧度量空间 \(K\) 到度量空间 \(Y\) 的连续映射一致连续。特别地，紧空间上的连续实函数有界，并取到最大值与最小值。
\end{proposition}
\begin{proof}
若一致连续性不成立，则存在 \(\varepsilon>0\) 与点对 \(x_n,y_n\in K\)，使 \(d_K(x_n,y_n)<1/n\)，但像点距离至少为 \(\varepsilon\)。取 \(x_n\) 的收敛子列，三角不等式说明对应的 \(y_n\) 趋于同一点，连续性给出矛盾。连续实函数的像是 \(\mathbb R\) 中的紧集，因此有界且闭；非空时它的上、下确界均属于像。
\end{proof}

这里的极值断言约定 \(K\ne\varnothing\)。它在几何测度论中常用来控制一个紧参数块上的导数，但不能据此给整个非紧区域一个统一常数。局部估计到整体估计之间，还需要覆盖或穷竭。

\subsection{紧邻域与可数穷竭}
\label{b1:subsec:A10203}

沿用\secref{b1:sec:1101}的约定，局部紧 Hausdorff 空间中的每一点有紧邻域。局部紧与全局紧不同，例如 \(\mathbb R^d\) 不是紧空间；局部紧也不保证可用可数个紧集覆盖，不可数离散空间便是反例。

\begin{lemma}\label{b1:lem:app-local-shrink}
在局部紧 Hausdorff 空间中，对 \(x\in U\)、\(U\) 开，可取开集 \(V\)，使 \(x\in V\subset\overline V\subset U\)，并且 \(\overline V\) 紧。若 \(K\subset U\) 紧，可取同样的 \(V\) 满足 \(K\subset V\subset\overline V\subset U\)。
\end{lemma}
\begin{proof}
取 \(x\) 的紧邻域 \(N\)。集合 \(F=N\setminus U\) 紧。用 Hausdorff 性逐点分开 \(x\) 与 \(F\)，再从 \(F\) 一侧取有限覆盖，可得互不相交的开集 \(O\ni x\)、\(W\supset F\)。于是 \(V=O\cap N^\circ\) 满足
\(\overline V\subset N\setminus W\subset U\)，且闭包作为 \(N\) 的闭子集紧。若 \(F\) 为空，直接取 \(O=X\)、\(W=\varnothing\)。紧集版本由有限个这样的 \(V\) 取并得到。
\end{proof}

这是\cref{b1:lem:lch-shrink}所用的缩小邻域步骤。这里把它放回拓扑背景，以便看清 Hausdorff 条件与紧性的不同任务。

\begin{proposition}\label{b1:prop:app-compact-exhaustion}
设 \(X\) 为局部紧 Hausdorff 空间，而且每个开覆盖都有可数子覆盖。则存在紧集列 \((K_n)\)，满足
\[
 K_n\subset K_{n+1}^{\circ},\quad X=\bigcup_{n=1}^{\infty}K_n^\circ.
\]
对任意紧集 \(K\subset X\)，有 \(K\subset K_n^\circ\) 对某个 \(n\) 成立。
\end{proposition}
\begin{proof}
以相对紧开集覆盖 \(X\)，再选可数子覆盖 \(V_1,V_2,\ldots\)。先选相对紧开集 \(W_1\) 包含 \(\overline V_1\)。已构造 \(W_n\) 后，用上一引理覆盖紧集 \(\overline W_n\cup\overline V_{n+1}\)，得到相对紧开集 \(W_{n+1}\) 包含它。令 \(K_n=\overline W_n\)，则 \(K_n\subset W_{n+1}\subset K_{n+1}^\circ\)，且这些内部覆盖 \(X\)。紧集 \(K\) 被有限个内部覆盖，因其递增，可取最大的一个。
\end{proof}

因此穷竭不是随便罗列一列紧集：内部包含关系保证任意固定紧集最终落在某一层内部，给截断、局部弱收敛与对角抽取留出余地。\secref{b1:sec:A103}将证明第二可数性提供这里所需的可数子覆盖条件。

% !TeX root = ../../Book1.tex
\section{可数性与乘积拓扑}
\label{b1:sec:A103}

可分性提供一列测试点，第二可数性提供一列测试开集。前者谈“点够不够多”，后者谈“邻域能否由可数种形状生成”。在度量空间中它们等价，在一般拓扑中却不能混用。中文教材中这两个名称可对照熊金城《点集拓扑讲义》第五版\cite[第5章]{Xiong2020Topology}；其积空间部分也适合作为本节的进一步阅读。

\subsection{从稠密点到可数基}
\label{b1:subsec:A10301}

\begin{definition}\label{b1:def:app-countability}
若 \(D\subset X\) 满足 \(\overline D=X\)，则称 \(D\) 在 \(X\) 中稠密；若存在可数稠密子集，则称 \(X\) \term[kefen]{可分}[separable]。若开集族 \(\mathcal G\) 满足每个开集均是其中若干成员的并，则称它为拓扑基。若有可数拓扑基，则称 \(X\) 满足\term[dier-keshuxing]{第二可数性}[second countability]。
\end{definition}

拓扑基的等价用法是：对 \(x\in U\)、\(U\) 开，可取 \(G\in\mathcal G\)，使 \(x\in G\subset U\)。稠密集则只要求与每个非空开集相交。两者之间少了对开集大小的控制，所以一般不能直接用稠密集的点替代一组基。

\begin{theorem}\label{b1:thm:app-separable-base}
第二可数空间可分，而且每个开覆盖都有可数子覆盖。对度量空间，可分性又推出第二可数性。第二可数空间的每个子空间仍第二可数；因此可分度量空间的每个子空间可分。
\end{theorem}
\begin{proof}
从可数基的每个非空成员中取一点，所得可数集与每个非空开集相交，故稠密。给定开覆盖，只考虑那些包含在某一覆盖成员中的基本开集，并为每个这样的基元选一个包含它的成员。这些选择至多可数，且覆盖全空间。

设 \((X,d)\) 有可数稠密集 \(D\)。以 \(D\) 中的点为中心、正有理数为半径的球组成可数族。给定 \(x\in U\)，先取 \(\varepsilon>0\)，使 \(B_d(x,\varepsilon)\subset U\)，再取 \(a\in D\) 满足 \(d(a,x)<\varepsilon/4\)，取有理数 \(r\) 满足 \(d(a,x)<r<\varepsilon/2\)。于是 \(x\in B_d(a,r)\subset U\)，这证明它是一组基。

若 \(\mathcal G\) 为 \(X\) 的可数基，则 \(\{A\cap G:G\in\mathcal G\}\) 是 \(A\) 的可数基。最后一句由此前两项应用于度量子空间得到。
\end{proof}

最后一句不能用 \(D\cap A\) 来证明：稠密集与子空间的交可能为空，例如 \(\mathbb Q\cap(\mathbb R\setminus\mathbb Q)\)。真正继承给子空间的是可数基，再由新基选择稠密点。

有限球网也提供可数点集。全有界空间对每个 \(n\) 取一个有限 \(1/n\)-网，把所有中心取并便得可数稠密集。因此紧度量空间可分；它的连续函数可常常通过可数点上的信息确定，但函数之间的一致估计仍需要连续性或等度连续性，不能仅凭“可数”二字省去。

\begin{proposition}\label{b1:prop:app-continuous-dense}
若 \(D\) 在 \(X\) 中稠密，\(f,g:X\rightarrow Y\) 连续，且 \(Y\) 为 Hausdorff 空间，则 \(f=g\) 只需在 \(D\) 上成立。若 \(X\) 可分，则连续像 \(f(X)\) 也可分。
\end{proposition}
\begin{proof}
若 \(f(x)\ne g(x)\)，用互不相交的开邻域分开这两个点，其反像的交给出 \(x\) 的开邻域，在其中 \(f\) 与 \(g\) 处处不同，与该邻域遇到 \(D\) 矛盾。第二句取可数稠密集 \(D\)；像空间中的非空相对开集有非空开反像，反像遇到 \(D\)，因此该开集遇到 \(f(D)\)。
\end{proof}

\subsection{有限坐标控制}
\label{b1:subsec:A10302}

\secref{b1:sec:0903}已经引入乘积拓扑并证明 Tychonoff 定理。这里补充其使用方法。设 \(X=\prod_{i\in I}X_i\)，坐标投影为 \(\pi_i\)。基本开集是
\[
 \bigcap_{i\in F}\pi_i^{-1}(U_i),\quad F\subset I\text{ 有限},\quad U_i\subset X_i\text{ 开}.
\]
不在 \(F\) 中的坐标完全不受限制。特别地，乘积拓扑是使全部坐标投影连续的最弱拓扑：任何这样的拓扑都必须包含上述有限交及其任意并。

\begin{proposition}\label{b1:prop:app-product-universal}
映射 \(f:Z\rightarrow\prod_iX_i\) 连续，当且仅当每个 \(\pi_i\circ f\) 连续。一个网在乘积拓扑中收敛，当且仅当它在每个坐标中收敛。
\end{proposition}
\begin{proof}
必要性由投影连续得到。充分性只需考察基本开集的反像：它是有限个坐标反像的交，因而开。对网的收敛，给定一个限制有限坐标的基本邻域，每个受限坐标都在某个指标以后满足要求。有向性保证这有限个指标有共同上界，于是从该上界起同时满足全部限制。
\end{proof}

若把“有限坐标”换成“所有坐标”，得到的更强拓扑一般不再有这一结论。例如 \(\mathbb R^{\mathbb N}\) 中的标准单位向量 \(e_n\) 逐坐标趋零，但在要求所有坐标同时属于 \((-1/2,1/2)\) 的盒状邻域中，没有一个 \(e_n\) 属于它。这个集合不是乘积拓扑的零点邻域。

\begin{theorem}[可数乘积的度量]\label{b1:thm:app-countable-product}
设 \((X_j,d_j)\)、\(j\ge1\)，是非空度量空间。下式给出其乘积拓扑的一个度量：
\[
 D(x,y)=\sum_{j=1}^{\infty}2^{-j}\min\{1,d_j(x_j,y_j)\}.
\]
若每个 \(X_j\) 完备，则 \(D\) 完备；若每个 \(X_j\) 可分，则乘积可分；若每个 \(X_j\) 紧，则乘积紧。
\end{theorem}
\begin{proof}
截断距离仍满足三角不等式，且上述级数绝对收敛，所以 \(D\) 是度量。给定限制有限坐标的基本邻域，对这些坐标取相应正距离阈值，再选足够小的 \(D\)-球，便落在其中。反之，给定 \(\varepsilon>0\)，先选 \(N\)，使 \(\sum_{j>N}2^{-j}<\varepsilon/2\)，再把前 \(N\) 个坐标限制在足够小的球中，使级数前段小于 \(\varepsilon/2\)。所得基本邻域包含于 \(D\)-球。故两种拓扑一致。

若 \((x^{(n)})\) 为 \(D\)-Cauchy 列，固定坐标 \(j\) 后，由其对应的一项小于总和，得到 \((x_j^{(n)})\) 是 \(d_j\)-Cauchy 列。逐坐标取极限 \(x_j\)，再用前段有限和与尾和估计，得 \(D(x^{(n)},x)\rightarrow0\)。

为证可分，对每个因子选可数稠密集 \(A_j\) 及基点 \(a_j\in A_j\)。考虑那些只有有限个坐标偏离 \(a_j\)，且每个坐标属于 \(A_j\) 的点。这个点集可数，并与每个非空基本开集相交。紧性可由\cref{b1:thm:tychonoff}直接得到；也可逐坐标抽收敛子列，再取对角子列，用上述尾和估计证明在 \(D\) 下收敛，最后应用度量紧性判据。
\end{proof}

这里可数个坐标使级数权重与对角子列都可用。一般不可数乘积的紧性仍成立，但不能从这一度量公式或可数次抽取推出，这正是 Tychonoff 定理在\chapref{b1:ch:09}中不可被简单对角法替代的原因。

\subsection{Borel 结构中的可数性}
\label{b1:subsec:A10303}

\begin{proposition}\label{b1:prop:app-product-borel}
若每个 \(X_j\) 第二可数，\(j\ge1\)，则
\[
 \mathcal B\left(\prod_{j=1}^{\infty}X_j\right)
 =\bigotimes_{j=1}^{\infty}\mathcal B(X_j),
\]
右侧指由各坐标投影的 Borel 反像生成的 \(\sigma\)-代数。
\end{proposition}
\begin{proof}
投影连续，故右侧包含于左侧。反向包含只需证明每个开集属于右侧。各因子选择可数基，由有限个坐标基元组成的基本开集只有可数多个，它们组成乘积的一组可数基。每个基本开集是有限个坐标 Borel 反像的交，每个开集是至多可数个这类集合的并，故属于右侧。
\end{proof}

证明中不是“开集是基本开集之并”就已足够，因为 \(\sigma\)-代数只对可数并封闭。可数基把任意并变成可数并，这是拓扑假设真正进入测度论的位置。\chapref{b1:ch:07}的乘积测度使用可测矩形，不能在一般空间中不经这一步就把乘积可测集与全部乘积 Borel 集等同。

% !TeX root = ../../Book1.tex
\section{闭集分离与连续延拓}
\label{b1:sec:A104}

在度量空间中，点到集合的距离天然给出连续函数。一般拓扑没有距离，分离闭集的开邻域便成为替代工具。\secref{b1:sec:1101}已经在局部紧空间中构造过紧支集截断；本节给出正规空间上的一般分离原理，并由一列逐次改进的近似建立连续延拓。

\subsection{正规性与邻域缩小}
\label{b1:subsec:A10401}

\begin{definition}\label{b1:def:app-normal}
若 \(X\) 的每个单点集闭，称它为 \(T_1\) 空间。若任意不同两点有互不相交的开邻域，称它为 Hausdorff 空间。若 \(X\) 为 \(T_1\) 空间，而且任意两个不交闭集均可分别包含在两个不交开集中，则称 \(X\) 为\term[zhenggui-kongjian]{正规空间}[normal space]。
\end{definition}

本书把 \(T_1\) 条件包含在正规空间的定义中；有些文献把它单列，读者比较分离公理时须先检查这一约定。以下实际的函数构造只用分开闭集的条件。

\begin{lemma}\label{b1:lem:app-normal-shrink}
正规空间中，若 \(F\subset U\)、\(F\) 闭且 \(U\) 开，则存在开集 \(V\)，使 \(F\subset V\subset\overline V\subset U\)。每个度量空间正规，每个紧 Hausdorff 空间也正规。
\end{lemma}
\begin{proof}
把不交闭集 \(F\) 与 \(X\setminus U\) 用不交开集 \(V,W\) 分开，则 \(\overline V\subset X\setminus W\subset U\)。

在度量空间中，单点闭。对非空不交闭集 \(A,B\)，函数 \(\mathrm{d}(x,A)=\inf_{a\in A}\mathrm{d}(x,a)\) 满足
\[
 |\mathrm{d}(x,A)-\mathrm{d}(y,A)|\le \mathrm{d}(x,y),
\]
因而连续，且零点集恰为 \(A\)。取
\[
 U=\{\,x:\mathrm{d}(x,A)<\mathrm{d}(x,B)/2\,\},\quad
 V=\{\,x:\mathrm{d}(x,B)<\mathrm{d}(x,A)/2\,\}.
\]
它们分别包含 \(A,B\)，且不相交。有空集时直接处理即可。

若 \(X\) 紧且 Hausdorff，先把一点 \(a\notin B\) 与闭集 \(B\) 分开：逐点选择不交邻域，在 \(B\) 一侧取有限覆盖，再把 \(a\) 一侧的邻域取交。对另一个紧闭集 \(A\) 中的点重复这一步，在 \(A\) 一侧再取有限覆盖，便得到分别包含 \(A,B\) 的不交开集。Hausdorff 性也保证 \(T_1\)。
\end{proof}

注意两个不交闭集之间的距离可以为零；上面的证明没有声称存在统一的正距离。例如平面中横轴与曲线 \(y=e^{-x}\) 是不交闭集，它们在无穷远处越来越近。距离函数在每个点的分离足以工作，不需要全局间隔。

\subsection{Urysohn 引理}
\label{b1:subsec:A10402}

\begin{theorem}[Urysohn]\label{b1:thm:app-urysohn}
设 \(X\) 正规，\(A,B\subset X\) 是不交闭集。存在连续函数 \(u:X\rightarrow[0,1]\)，使 \(u=0\) 于 \(A\)，\(u=1\) 于 \(B\)。
\end{theorem}
\begin{proof}
若有一个集合为空，取相应常函数。以下设二者非空。令 \(U_1=X\setminus B\)，利用缩小引理取开集 \(U_0\)，使 \(A\subset U_0\subset\overline U_0\subset U_1\)。在 \(0,1\) 之间插入 \(1/2\)，令
\(\overline U_0\subset U_{1/2}\subset\overline U_{1/2}\subset U_1\)。
继续在每两个相邻的已选二进有理数之间插入中点，并对闭集 \(\overline U_r\) 与开集 \(U_s\) 应用缩小引理。这样对所有二进有理数
\(\mathcal D=\{\,k2^{-n}:0\le k\le2^n,\ n\ge0\,\}\)
得到开集 \(U_r\)，满足
\[
 \overline U_r\subset U_s\quad(0\le r<s\le1).
\]
定义
\[
 u(x)=\inf\bigl(\{\,r\in\mathcal D:x\in U_r\,\}\cup\{1\}\bigr).
\]
若 \(x\in A\)，则 \(x\in U_0\)，故 \(u(x)=0\)；若 \(x\in B\)，它不属于任何 \(U_r\)，故 \(u(x)=1\)。

对 \(0<t\le1\)，有
\[
 \{u<t\}=\bigcup_{\substack{r\in\mathcal D\\r<t}}U_r.
\]
对 \(0\le t<1\)，还有
\[
 \{u>t\}
 =\bigcup_{\substack{r\in\mathcal D\\t<r<1}}
       \bigl(X\setminus\overline U_r\bigr).
\]
为核对后一等式，若 \(x\notin\overline U_r\)，则所有 \(s<r\) 都不可能满足 \(x\in U_s\)，于是 \(u(x)\ge r>t\)。反之，若 \(u(x)>t\)，取 \(t<r<s<u(x)\) 的二进有理数。此时 \(x\notin U_s\)，由 \(\overline U_r\subset U_s\) 得 \(x\notin\overline U_r\)。两类逆像均开，端点之外的逆像又是空集或全空间，故 \(u\) 连续。
\end{proof}

关键不是只让 \(U_r\subset U_s\)，而是把闭包也放进去。这个余量同时控制严格下水平集与严格上水平集，从而得到连续性。Fremlin《Measure Theory》\cite[Appendix to Volume 4, 4A2F(d)]{Fremlin2016Measure}列有这一引理和下面的延拓结论，适合查阅不同分离公理下的版本。

\ProseParagraphBreak

若只在度量空间中分开两个非空闭集，可直接写
\[
 u(x)=\frac{\mathrm{d}(x,A)}{\mathrm{d}(x,A)+\mathrm{d}(x,B)}.
\]
分母在每个点为正，因为闭集 \(A,B\) 不交。若其中一个集合为空，定理所需函数可直接取常数：\(A=\varnothing\) 时取 \(u\equiv1\)，\(B=\varnothing\) 时取 \(u\equiv0\)；二者皆空时任选其一。一般正规空间上的二进插值，是在没有距离时构造同类函数的方法。

\subsection{Tietze 延拓及值域控制}
\label{b1:subsec:A10403}

下面的结果称为\term[tietze-yantuo]{Tietze 延拓定理}[Tietze extension theorem]。%
\footnote{熊金城《点集拓扑讲义》第五版\cite[第6.3节]{Xiong2020Topology}使用“Tie\-tze 扩张定理”这一名称，故同时登记“Tietze 扩张定理”\termsee[tietze-kuozhang]{Tietze 扩张定理}{Tietze 延拓定理}供检索。此处“扩张”与“延拓”指同一闭集上的连续函数延伸问题。}

\begin{theorem}[Tietze]\label{b1:thm:app-tietze}
设 \(X\) 正规，\(F\subset X\) 闭。每个连续实函数 \(f:F\rightarrow\mathbb R\) 都有连续延拓 \(g:X\rightarrow\mathbb R\)。若 \(f\) 的值域包含于闭区间 \([a,b]\)，可让 \(g\) 也取值于该区间。
\end{theorem}
\begin{proof}
先证明一个逼近步骤。若 \(h:F\rightarrow[-M,M]\) 连续，\(M>0\)，则闭集
\[
 A=\{\,x\in F:h(x)\le-M/3\,\},\quad
 B=\{\,x\in F:h(x)\ge M/3\,\}
\]
也是 \(X\) 中的闭集，且互不相交。用 Urysohn 引理构造 \(v:X\rightarrow[-M/3,M/3]\)，使 \(v=-M/3\) 于 \(A\)，\(v=M/3\) 于 \(B\)。在 \(F\) 上逐段比较可得 \(|h-v|\le2M/3\)：两端区间利用规定的端点值，中间区间利用 \(h,v\) 均落在 \([-M/3,M/3]\)。

设 \(|f|\le M\)。以 \(h_0=f\) 开始，递归把上述步骤用于当前余项，得到连续函数 \(v_n:X\rightarrow\mathbb R\)、\(n\ge0\)，使
\[
 \|v_n\|_\infty\le\frac M3\left(\frac23\right)^n,\quad
 \left|f-\sum_{j=0}^{n}v_j\right|
 \le M\left(\frac23\right)^{n+1}\quad\hbox{于 }F.
\]
级数 \(g=\sum_{n\ge0}v_n\) 一致收敛，故连续；余项估计保证 \(g=f\) 于 \(F\)，几何级数又给出 \(|g|\le M\)。平移、缩放便处理一般闭区间；退化区间取常函数。

对可能无界的 \(f\)，令 \(h=f/(1+|f|)\)，这是 \(F\) 上取值于 \((-1,1)\) 的连续函数。由有界情形，先延拓为 \(H:X\rightarrow[-1,1]\)。集合 \(E=\{\,x:|H(x)|=1\,\}\) 闭且不交 \(F\)。再用 Urysohn 引理取 \(\eta:X\rightarrow[0,1]\)，使 \(\eta=1\) 于 \(F\)、\(\eta=0\) 于 \(E\)。令 \(G=\eta H\)。在 \(E\) 上 \(G=0\)，在其外 \(|G|<1\)，所以
\[
 g=\frac{G}{1-|G|}
\]
是整个 \(X\) 上有限的连续实函数，且在 \(F\) 上等于 \(f\)。
\end{proof}

无界情形不能把一个 \([-1,1]\) 值延拓直接代入逆变换：它可能在 \(F\) 之外取到端点。第二次分离恰好把这些端点消掉，同时不改变 \(F\) 上的值。

\ProseParagraphBreak

这里也没有断言延拓唯一，或给出一个线性的延拓选择。比如从闭集 \(\{0\}\subset\mathbb R\) 延拓零函数，有无穷多个选择。与\cref{b1:prop:app-uniform-extension}相比，一个在闭集上，另一个在稠密集上；闭性提供“可以自由补函数”的空间，稠密性则提供唯一性。

\begin{corollary}\label{b1:cor:app-compact-extension}
设 \(X\) 为局部紧 Hausdorff 空间，\(K\subset X\) 紧，且 \(U\supset K\) 开。每个连续实函数 \(f:K\rightarrow\mathbb R\) 都可延拓为 \(g\in C_c(X)\)，满足 \(\operatorname{supp}g\subset U\) 以及 \(\|g\|_\infty\le\|f\|_\infty\)。
\end{corollary}
\begin{proof}
选相对紧开集 \(V\)，使 \(K\subset V\subset\overline V\subset U\)。紧 Hausdorff 空间 \(\overline V\) 正规；由 Tietze 定理把 \(f\) 延拓到 \(\overline V\)，不增加上确界范数。再由\cref{b1:lem:lch-cutoff}取 \(0\le\chi\le1\)，使 \(\chi=1\) 于 \(K\) 且 \(\operatorname{supp}\chi\subset V\)。在 \(V\) 内令 \(g=\chi f_{\mathrm{ext}}\)，在 \(V\) 外补零。由于 \(\chi\) 的支集是 \(V\) 内的紧集，拼接处函数在邻域中恒为零，故延拓连续，并有要求的支集及范数控制。
\end{proof}

此推论只用一个紧邻域的正规性，不要求整个局部紧 Hausdorff 空间正规。这也是在一般 Radon 测度论中宜先把问题局部化的原因。

% !TeX root = ../../Book1.tex
\section{波兰空间与相容的完备度量}
\label{b1:sec:A105}

第\ref{b1:sec:1204}节的波兰空间结合两种优势：可数基使测试条件可以可数化，某个相容度量的完备性使 Cauchy 逼近留在空间内。这里补充这一类空间在构造下的稳定性。由于定义中的完备度量只要求存在，我们可以改变距离的数值，只要不改变开集。

\subsection{开子集为什么仍可完备度量化}
\label{b1:subsec:A10501}

\begin{proposition}\label{b1:prop:app-polish-operations}
波兰空间的闭子空间、可数乘积仍是波兰空间。若 \(X\) 为波兰空间且 \(G\subset X\) 是可数个开集之交，则 \(G\) 也是波兰空间；特别地，开子空间也是。
\end{proposition}
\begin{proof}
闭子空间的断言由完备性继承与可分度量子空间的可分性得到。可数乘积用\cref{b1:thm:app-countable-product}，对每个因子先选相容完备度量。

令 \(G=\bigcap_{j\ge1}U_j\)，在 \(X\) 上固定相容完备度量 \(d\)。若 \(G\) 为空，断言直接成立。对非空补集 \(X\setminus U_j\)，在 \(G\) 上设
\[
 h_j(x)=\frac1{d(x,X\setminus U_j)};
\]
若补集为空，令 \(h_j=0\)。每个分母在 \(G\) 上为正，故 \(h_j\) 连续。考虑图映射
\[
 \Phi:G\longrightarrow X\times\mathbb R^{\mathbb N},
 \quad x\longmapsto\bigl(x,h_1(x),h_2(x),\ldots\bigr).
\]
它与第一坐标投影互逆，因此是到其像的同胚。下面证明像闭。

设图上序列 \(\Phi(x_n)\) 在乘积中趋向 \((x,t_1,t_2,\ldots)\)。若 \(x\notin U_j\) 且该补集非空，则 \(d(x_n,X\setminus U_j)\le d(x_n,x)\rightarrow0\)，从而 \(h_j(x_n)\rightarrow+\infty\)，与 \(t_j\) 有限矛盾。所以 \(x\in G\)，各坐标连续性再给出 \(t_j=h_j(x)\)。乘积可度量，序列判据足以说明像闭。该乘积有相容完备度量，闭子空间因而完备；把它拉回 \(G\)，便得与原子空间拓扑相容的完备度量。可分性已由 \(G\subset X\) 继承。
\end{proof}

这个构造并没有把边界点补进 \(G\)。恰恰相反，倒数距离使趋向被删边界的序列在某个新坐标上逃向无穷，从而不再是新度量下的 Cauchy 列。可把拉回度量具体写成
\[
 \rho(x,y)=\min\{1,d(x,y)\}
 +\sum_{j=1}^{\infty}2^{-j}\min\{1,|h_j(x)-h_j(y)|\}.
\]
图闭的论证说明它完备，而有限前段加小尾和的论证说明它不改变拓扑。这一形式与 Fremlin《Measure Theory》卷四\cite[4A2M、4A2Q]{Fremlin2016Measure}中的完备度量构造相联系。

\subsection{例子与不能省去的条件}
\label{b1:subsec:A10502}

欧氏空间的开集、闭集都是波兰空间。无理数集也是，因为它是可数个开集 \(\mathbb R\setminus\{q\}\)、\(q\in\mathbb Q\)，的交。无理数集的通常距离并不完备，例如无理数列可以趋向有理数；上一构造给出的是另一相容完备度量。

有理数集却不是波兰空间。若其通常拓扑可由某个完备度量给出，则\cref{b1:thm:baire}适用。每个有理数单点闭且没有内点，所以 \(\mathbb Q\) 是可数个无处稠密闭集的并，与非空完备度量空间的 Baire 性矛盾。这里无处稠密性只依赖拓扑，因此换成任何相容度量也不能消除矛盾。

由此还得到一个值得保留的区别：波兰空间的任意子集不一定波兰，即使它是 Borel 集。\(\mathbb Q\) 是 \(\mathbb R\) 中的可数并闭集，却不波兰。相反，可数交开集的结论有上面的明确构造作为保证，不能把“Borel”三个字直接替换进去。

\begin{proposition}\label{b1:prop:app-polish-functions}
若 \(K\) 是紧度量空间，则实连续函数空间 \(C(K)\) 在一致范数下为可分 Banach 空间，因而是波兰空间。
\end{proposition}
\begin{proof}
一致 Cauchy 列逐点有极限，并由共同的 Cauchy 估计一致收敛；一致极限连续，所以 \(C(K)\) 完备。以下证可分，空 \(K\) 的情形直接成立。

固定可数稠密集 \(D\subset K\)。考虑所有有限个 \(a_i\in D\)、正有理数 \(r\)、有理数 \(q_i\) 所给出的函数
\[
 g(x)=\frac{\sum_{i=1}^N q_i\max\{0,r-d(x,a_i)\}}
 {\sum_{i=1}^N\max\{0,r-d(x,a_i)\}},
\]
只保留分母在整个 \(K\) 上为正的选择。它们组成可数族。给定 \(f\in C(K)\) 与 \(\varepsilon>0\)，一致连续性允许选正有理数 \(r\)，使 \(d(x,y)<r\) 时 \(|f(x)-f(y)|<\varepsilon/2\)。从 \(D\) 选有限个中心，其 \(r/2\)-球覆盖 \(K\)，再选有理数 \(q_i\) 满足 \(|q_i-f(a_i)|<\varepsilon/2\)。上述分母为正，而每个非零权重都满足 \(d(x,a_i)<r\)。因此 \(g(x)\) 是若干与 \(f(x)\) 相差小于 \(\varepsilon\) 的数的凸组合，故 \(\|g-f\|_\infty<\varepsilon\)。
\end{proof}

这解释了为什么连续路径空间可以纳入第十二章的测度收敛理论。例如 \(C([0,1])\) 完备且可分，但它的闭单位球不紧；不同频率或不同窄峰的函数会破坏等度连续性。波兰性给出适宜的状态空间，紧族条件仍需另行证明。

\subsection{返回测度与局部分析}
\label{b1:subsec:A10503}

本附录的工具可以沿两条方向使用。若研究紧支集测试与局部积分，局部紧性提供紧邻域，闭集分离提供连续截断，可数基则进一步提供紧穷竭。若研究概率分布或随机路径，可分完备度量提供统一的可数逼近与序列极限，未必存在紧邻域，因此不能把 \(C_c\) 的方法原样搬过去。

在第\ref{b1:sec:1204}节证明单个概率测度的紧性时，可分性给出可数小球覆盖，再在每一尺度选有限子族；把各尺度的覆盖同时满足，得到全有界的闭集，完备性最后使它紧。现在读者可以清楚辨认这三个步骤：可数化、有限逼近、保留极限。它们不是同一种“紧性”的不同叫法，而是构造紧集时先后承担的不同任务。


\begin{chaptersummary}
拓扑记录邻域和连续性，完备性则还记录 Cauchy 逼近所使用的距离。紧度量空间恰好是完备而全有界的空间；可数基使开覆盖和 Borel 生成过程可以可数化。正规性把闭集分离转化为连续函数，并由几何递减的误差构造闭集上的连续延拓。波兰空间允许选择合适的完备度量，但并不自动局部紧；使用紧支集方法和使用概率测度紧族方法时，应分别检查所需假设。
\end{chaptersummary}

\BookExercises

% 以下八题均自拟；用于辨析假设、构造拓扑反例和连接函数空间，不改编未核实的原题。
\begin{exercise}\label{b1:ex:A1-contraction}
在 \(X=\mathbb Q\cap[1,2]\) 上取通常距离，令 \(T(x)=(x+2)/(x+1)\)。证明 \(T(X)\subset X\)，且 \(T\) 是常数不超过 \(1/4\) 的压缩映射，但在 \(X\) 中没有不动点。说明其迭代为何构成 Cauchy 列，并在完备化中找到极限。
\end{exercise}
\begin{solution}
函数 \(T\) 在 \([1,2]\) 上取值于 \([4/3,3/2]\)，并保持有理数。直接相减可得
\[
 |T(x)-T(y)|=\frac{|x-y|}{(x+1)(y+1)}\le\frac14|x-y|.
\]
不动点方程为 \(x^2=2\)，故在 \(X\) 中无解。令 \(x_{n+1}=T(x_n)\)，则相邻差被 \(4^{-n}|x_1-x_0|\) 控制，尾部差以几何级数估计，故是 Cauchy 列。\(X\) 的完备化为 \([1,2]\)；连续延拓的极限满足同一方程，因而为 \(\sqrt2\)。这说明不动点论证需要完备性，不能只检查压缩常数。
\end{solution}

\begin{exercise}\label{b1:ex:A1-l2-local}
证明 \(\ell^2\) 的闭单位球有界且完备，但不紧；进一步证明 \(\ell^2\) 不局部紧。另一方面，说明它仍是波兰空间。
\end{exercise}
\begin{solution}
空间 \(\ell^2\) 完备，闭单位球也完备。标准单位向量两两距离为 \(\sqrt2\)，因此球不全有界，不能紧。若零点有紧邻域，该邻域包含半径 \(r>0\) 的开球，便包含两两距离为 \(r/\sqrt2\) 的点 \((r/2)e_n\)，与紧集全有界矛盾。平移说明每一点均无紧邻域。有限支集有理数序列在实 \(\ell^2\) 中可数且稠密；复情形改用有理实部和虚部。因此它完备、可分，是波兰空间。
\end{solution}

\begin{exercise}\label{b1:ex:A1-product-sequences}
设 \(I\) 不可数，在 \(X=\{0,1\}^I\) 上取乘积拓扑。令 \(\Sigma\) 为所有至多在可数个坐标取值 \(1\) 的点。证明 \(\Sigma\) 稠密且对序列极限封闭，却不是闭集。说明为何不能用序列检测这个空间中的闭性。
\end{exercise}
\begin{solution}
每个非空基本开集只限制有限个坐标；在其余坐标取零便给出 \(\Sigma\) 中的一点，所以 \(\Sigma\) 稠密。若 \(x_n\in\Sigma\) 且 \(x_n\rightarrow x\)，各 \(x_n\) 的非零坐标集之并 \(J\) 可数。对 \(i\notin J\)，全部 \(x_n(i)=0\)，逐坐标收敛给出 \(x(i)=0\)，故 \(x\in\Sigma\)。但恒等于 \(1\) 的点不属于 \(\Sigma\)，所以它是稠密真子集而不闭。这个例子不与度量空间的序列判据矛盾，反而说明此不可数乘积不能具有给出其拓扑的度量。
\end{solution}

\begin{exercise}\label{b1:ex:A1-distance-test}
设 \(Y\) 为非空可分度量空间，\(D\subset Y\) 为可数稠密集，\((S,\mathcal A)\) 为可测空间。证明 \(f:S\rightarrow Y\) 是 Borel 可测的，当且仅当对每个 \(a\in D\)，实函数 \(s\mapsto d(f(s),a)\) 可测。
\end{exercise}
\begin{solution}
必要性由距离函数连续得到。反之，假设这些实函数可测，则所有以 \(D\) 中的点为中心、正有理半径的球的反像可测。它们是可数拓扑基；任意开集为其中至多可数个成员的并，故其反像可测。再由 Borel \(\sigma\)-代数的生成性质得到 \(f\) 可测。可数性用在开集到可测并的步骤，不是用于任意不可数并也可测的错误推断。
\end{solution}

\begin{exercise}\label{b1:ex:A1-vector-extension}
设 \(X\) 正规，\(F\subset X\) 闭，\(f:F\rightarrow\mathbb R^m\) 连续且 \(|f|\le1\)。证明存在连续延拓 \(g:X\rightarrow\mathbb R^m\)，仍满足 \(|g|\le1\)。为什么逐坐标延拓后还需要一步处理？
\end{exercise}
\begin{solution}
由 Tietze 定理逐坐标延拓，得连续 \(G:X\rightarrow\mathbb R^m\)。即使每个坐标都限制在 \([-1,1]\)，也只能直接保证 \(|G|\le\sqrt m\)，而不是单位球值域。定义连续映射
\[
 R(v)=\frac{v}{\max\{1,|v|\}},\quad g=R\circ G.
\]
它在闭单位球上为恒等，故 \(g=f\) 于 \(F\)，并且全域取值于该球。复值函数可视为 \(m=2\) 的情形，同样保持圆盘而不只保持外接正方形。
\end{solution}

\begin{exercise}\label{b1:ex:A1-hyperspace}
设 \((X,d)\) 为非空紧度量空间，以 \(\mathcal K(X)\) 表示全部非空紧子集。对 \(A,B\in\mathcal K(X)\) 规定
\[
 d_H(A,B)=\max\left\{\sup_{a\in A}d(a,B),\ \sup_{b\in B}d(b,A)\right\}.
\]
证明 \(d_H\) 为度量，并证明 \(\mathcal K(X)\) 在此度量下紧。
\end{exercise}
\begin{solution}
各上确界有限；若 \(d_H(A,B)=0\)，闭性说明 \(A\subset B\) 且 \(B\subset A\)。由
\(d(a,C)\le d(a,B)+\sup_{b\in B}d(b,C)\)
及反向同类估计，取上确界得到三角不等式。

给定 \(X\) 的有限 \(\varepsilon\)-网 \(\{x_1,\ldots,x_N\}\)。对每个非空紧集 \(A\)，保留那些 \(\varepsilon\)-球与 \(A\) 相交的网点，所得非空有限集 \(F_A\) 满足 \(d_H(A,F_A)\le\varepsilon\)。这样的有限集只有有限多个，故 \(\mathcal K(X)\) 全有界。

设 \((A_n)\) 是 \(d_H\)-Cauchy 列。令
\[
 A=\bigcap_{N=1}^{\infty}\overline{\bigcup_{n\ge N}A_n}.
\]
右侧是紧空间 \(X\) 中递减非空闭集的交，所以 \(A\) 非空且紧。给定 \(\varepsilon>0\)，选 \(N\)，使 \(n,m\ge N\) 时 \(d_H(A_n,A_m)<\varepsilon\)。固定 \(n\ge N\) 与 \(a\in A_n\)，在每个 \(A_m\)、\(m\ge N\)，选 \(a_m\) 距 \(a\) 不超过 \(\varepsilon\)。取趋于无穷的指标的收敛子列，极限属于 \(A\)，所以 \(d(a,A)\le\varepsilon\)。反过来，每个 \(b\in A\) 可由 \(b_j\in A_{m_j}\)、\(m_j\to\infty\)，逼近；故 \(d(b,A_n)\le\varepsilon\)。于是 \(d_H(A_n,A)\le\varepsilon\)。这证明完备，结合全有界性得到紧性。
\end{solution}

\begin{exercise}\label{b1:ex:A1-polish-image}
给出一个波兰空间及其连续满射，使像空间在给定拓扑下不是波兰空间。另证明无理数集不是实数直线中的可数并闭集。
\end{exercise}
\begin{solution}
取带离散距离的 \(\mathbb N\)，它完备且可分。把自然数按某个枚举映到带通常拓扑的 \(\mathbb Q\)，任何这样的映射都连续，因为离散空间的任意子集开；像 \(\mathbb Q\) 不是波兰空间。

若无理数集是 \(\mathbb R\) 中可数并闭集，则 \(\mathbb Q\) 是可数交开集，按\cref{b1:prop:app-polish-operations}应为波兰空间，矛盾。这并不否定连续像的可分性，而是表明可完备度量化性不会被一般连续满射保留。
\end{solution}

\begin{exercise}\label{b1:ex:A1-local-functions}
在 \(C(\mathbb R^d)\) 上令 \(K_n=\{\,x:|x|\le n\,\}\)，并设
\[
 \rho(f,g)=\sum_{n=1}^{\infty}2^{-n}\min\{1,\|f-g\|_{C(K_n)}\}.
\]
证明此度量的收敛恰为在每个紧集上一致收敛，并证明它完备且可分。
\end{exercise}
\begin{solution}
每个紧集包含于某个 \(K_n\)，因此在所有 \(K_n\) 上的一致收敛等价于在每个紧集上一致收敛。级数的有限前段与尾和估计又说明它等价于 \(\rho\)-收敛。

把函数送到其限制族，得到
\[
 C(\mathbb R^d)\longrightarrow\prod_{n\ge1}C(K_n).
\]
像由闭条件 \(f_{n+1}|_{K_n}=f_n\) 给出。满足这些条件的连续函数族能拼成全域连续函数，因为 \(K_n^\circ\) 覆盖整个空间；所以像正好是这个闭子空间。各 \(C(K_n)\) 在一致范数下完备且可分，可数乘积的标准截断度量拉回恰为 \(\rho\)。闭子空间完备，可分度量子空间可分，由此得到两项结论。这里得到的是紧集上一致收敛的完备距离，不是要求函数全域有界的一致范数。
\end{solution}
